% Generated draft: requires mathematical and source review.
\documentclass[11pt]{article}
\usepackage[margin=25mm]{geometry}
\usepackage{fontspec}
\setmainfont{Latin Modern Roman}
\usepackage{amsmath,amssymb,mathtools,unicode-math}
\setmathfont{Latin Modern Math}
\usepackage{enumitem,needspace,xurl,hyperref}
\setlength{\parindent}{0pt}
\setlength{\parskip}{5pt}
\setlength{\emergencystretch}{3em}
\newcommand{\N}{\mathbb{N}}
\newcommand{\Z}{\mathbb{Z}}
\newcommand{\Q}{\mathbb{Q}}
\newcommand{\R}{\mathbb{R}}
\newcommand{\C}{\mathbb{C}}
\newcommand{\F}{\mathbb{F}}
\newcommand{\A}{\mathbb{A}}
\newcommand{\PP}{\mathbb{P}}
\newcommand{\E}{\mathbb{E}}
\newcommand{\eps}{\varepsilon}
\newcommand{\dd}{\,\mathrm{d}}
\begin{document}
\section*{UnsolvedMath ID: 30000552}
\textbf{Source problem code:} catalogue entry
\subsection{Definitions and mathematical statement}
Let \(M\) carry two metrics \(d_1,d_2\), each making it a complete length space. Complete means every Cauchy sequence converges. Length means that \(d_i(x,y)\) is the infimum of the \(d_i\)-lengths of continuous paths joining \(x\) to \(y\); it does not by itself assert that a minimizing geodesic exists. Let the group \(\Gamma\) have the intended action on the common set \(M\), with
\[
 d_i(\gamma x,\gamma y)=d_i(x,y)
 \quad(i=1,2;\ \gamma\in\Gamma;\ x,y\in M),
\]
and assume the action is cocompact for both metrics. In the standard quotient-topology meaning, each quotient is compact. The source does not add properness of the action or of either metric space.

The asymptotic hypothesis is
\[
 \frac{d_1(x,y)}{d_2(x,y)}\longrightarrow1
 \quad\text{as }d_1(x,y)\to\infty.
\]
Its usual all-pairs interpretation is
\[
 \forall\varepsilon>0\ \exists R>0\ \forall x,y\in M,
 \quad d_1(x,y)\geq R\ \Longrightarrow\
 \left|\frac{d_1(x,y)}{d_2(x,y)}-1\right|<\varepsilon.
\]
The positive cutoff excludes \(x=y\), so the quotient is defined. Equivalently, the ratio must tend to one along every sequence of pairs with \(d_1\)-distance tending to infinity, not just pairs on a chosen ray or with one fixed endpoint.

The question asks whether these hypotheses imply
\[
 \exists C<\infty\ \forall x,y\in M,\qquad
 |d_1(x,y)-d_2(x,y)|\leq C.
\]
The constant can depend on the two metrics and the action but not on the pair. This is an additive bounded-difference conclusion on the given set, without changing points by a quasi-isometry or passing to a group orbit. Relative asymptotic agreement is the stated assumption; an additive bound is not built into it.

The source does not explicitly require the metrics to induce the same topology, specify whether one common compact fundamental set exists, or assume \(d_1\) unbounded. If there are no pairs with \(d_1\to\infty\), the asymptotic condition is vacuous and the remaining hypotheses must still be interpreted as printed. No local compactness, geodesicity, proper action or unboundedness condition is inserted to repair that possible scope issue.

The literature triage checked 2026-08-21 reports no authoritative resolution of the extracted claim located in its targeted search. This open label and the editorial verification note remain dated metadata, not independently verified current conclusions.
\subsection{Short English statement}
If two complete length metrics admit the same cocompact isometric group action and their ratios tend uniformly to one at large distances, must their difference be bounded on every pair of points? The additive conclusion is stronger than relative agreement. Topology compatibility, properness and unboundedness are not supplied as extra hypotheses.
\subsection{Sources}
\begin{itemize}
\item Local UnsolvedMath registry, permanent ID 30000552, source code catalogue entry, ``Asymptotically Equivalent Cocompact Metrics,'' including original and cleaned all-pairs question and verification metadata.
\item \emph{Geometric Group Theory, Hyperbolic Dynamics and Symplectic Geometry (2007)}, Oberwolfach Reports, pp.\ 1991--2058, DOI 10.4171/owr/2006/33: \url{https://doi.org/10.4171/owr/2006/33}. This is also the sole supplied literature reference. The citation uses 2007, while the DOI and proposed-year metadata use 2006.
\item Literature triage checked 2026-08-21; local record updated 2026-09-16. No external source has been visited or independently checked here.
\end{itemize}
\end{document}
